\documentclass[11pt]{article}

\usepackage[utf8]{inputenc}
\usepackage[T1]{fontenc}
\usepackage{lmodern}
\usepackage[a4paper,margin=1in]{geometry}
\usepackage{amsmath,amssymb,amsthm}
\usepackage{graphicx}
\usepackage{booktabs}
\usepackage{array}
\usepackage{xcolor}
\usepackage{microtype}
\usepackage[hidelinks]{hyperref}
\usepackage{xr-hyper}
\usepackage[numbers,sort&compress]{natbib}

% Cross-reference the main manuscript's equation/section numbers with an M- prefix.
\externaldocument[M-]{main}

\newcommand{\Gfac}{\textit{G}}
\newcommand{\info}{\mathcal{I}}
\newcommand{\unitinfo}{\mathcal{J}}
\newcommand{\eED}{\mathbb{E}}

% Supplementary numbering: S-prefixed sections, equations, theorem-like blocks.
\renewcommand{\thesection}{S\arabic{section}}
\renewcommand{\thesubsection}{S\arabic{section}.\arabic{subsection}}
\renewcommand{\theequation}{S\arabic{equation}}
\renewcommand{\thetable}{S\arabic{table}}
\renewcommand{\thefigure}{S\arabic{figure}}

\theoremstyle{plain}
\newtheorem{proposition}{Proposition}
\newtheorem{lemma}{Lemma}
\newtheorem{corollary}{Corollary}
\theoremstyle{definition}
\newtheorem{definition}{Definition}

% Main-text cross-reference helper: \MEq{label} -> "Eq. (M-3)" style.
\newcommand{\MEq}[1]{Eq.~(\ref{M-#1})}
\newcommand{\MSec}[1]{\S\ref{M-#1}}
\newcommand{\MFig}[1]{Fig.~\ref{M-#1}}

\title{\bfseries Supplementary Material\\[4pt]
\large A transdiagnostic dimensional map as a self-calibrating measurement
instrument: projection, information, and value of information under structured
missingness}

\author{Andrii Kulakovskyi \and Ophélia Godin \and Léah Camarcat \and
Engin Altunlu \and Marion Leboyer}

\date{\today}

\begin{document}
\maketitle

\noindent
This document gives the full derivations that the Methods section of the main text
(\MSec{sec:m-model}--\MSec{sec:m-calib}) states as results and defers here. It is
self-contained: every object it uses is defined below, and cross-references of the
form ``\MEq{eq:m-scores}'' point to the numbered equations of the main manuscript.
Sections~\ref{sec:model} and~\ref{sec:project} develop the measurement model and the
projection machinery from first principles; Sections~\ref{sec:slope}--\ref{sec:calib}
derive in full the five instrument properties that are the paper's contribution --- the
loading as a measurement slope, the information accounting and per-family Fisher
information, adaptive sampling and value of information, the archetypal summary and its
reconstruction fidelity, and the calibration protocol. Propositions are proved;
everything reported in the main text is an evaluation of the objects defined here.

\tableofcontents
\bigskip
\hrule
\bigskip

%=============================================================================
\section{The measurement model from first principles}
\label{sec:model}

\subsection{Generative factor model and local independence}
\label{sub:gen}

Index patients $i=1,\dots,N$, indicators $j=1,\dots,J$, specific factors
$k=1,\dots,K$ and one general factor \Gfac. Each patient carries a factor vector
$f_i=(G_i,D_{i1},\dots,D_{iK})^{\top}\in\mathbb{R}^{K+1}$ with
\begin{equation}
  f_i \sim \mathcal{N}_{K+1}\!\left(\mathbf{0},\,\Phi\right),
  \qquad \operatorname{diag}(\Phi)=\mathbf{1},
  \label{eq:factors}
\end{equation}
the factor \emph{correlation} matrix $\Phi$ carrying a unit diagonal to fix the latent
scale. Each indicator is a linear reflection of the factors plus item-specific noise
through a single linear predictor,
\begin{equation}
  \eta_{ij} = \alpha_j + \lambda_{jG}\,G_i + \sum_{k=1}^{K}\lambda_{jk}\,D_{ik}
  + \beta_j^{\top} c_i ,
  \label{eq:linpred}
\end{equation}
with item intercept $\alpha_j$, general loading $\lambda_{jG}$, specific loadings
$\lambda_{jk}$ (row $j$ of the loading matrix
$\Lambda=[\,\lambda_G\mid\Lambda_D\,]\in\mathbb{R}^{J\times(K+1)}$), and item-level
covariates $c_i$ that adjust an indicator's mean but are never factor indicators. This
is \MEq{eq:measure} of the main text. Stacking one patient's items,
\begin{equation}
  x_i=\mu_i+\lambda_G G_i+\Lambda_D D_i+\varepsilon_i,
  \qquad \mu_i=\alpha+B c_i,\qquad \varepsilon_i\sim\mathcal{N}(\mathbf 0,\Psi),
  \label{eq:xstack}
\end{equation}
with diagonal residual covariance $\Psi=\operatorname{diag}(\sigma_1^2,\dots,\sigma_J^2)$.

\paragraph{Local independence.} Given a patient's latent scores the items are
conditionally independent,
\begin{equation}
  p(x_i\mid f_i,\Theta)=\prod_{j\in O_i} p(x_{ij}\mid f_i,\Theta),
  \label{eq:localindep}
\end{equation}
i.e.\ the factors are the only reason items covary. In the Gaussian linear model this is
identical to ``$\Psi$ is diagonal'': local independence and diagonal residual covariance
are two faces of one assumption. It is also why filling missing cells is forbidden --- an
invented value injects covariance the model would mistake for a factor. Because $f_i$ is
unobserved we integrate it out and multiply across independent patients,
\begin{equation}
  L(\Theta)=\prod_{i=1}^{N}\int\Bigl[\prod_{j\in O_i} p(x_{ij}\mid f_i,\Theta)\Bigr]\,
  \mathcal{N}(f_i;\mathbf 0,\Phi)\,\mathrm{d}f_i .
  \label{eq:fulllik}
\end{equation}
Equation~\eqref{eq:linpred} is a bifactor exploratory structural-equation model: every
indicator may load on \Gfac{} and on its home specific factor, while off-home cells are
governed by a sparsity-inducing prior (\S\ref{sub:rotation}) that keeps the map mostly
simple-structure yet lets a few evidence-backed cross-loadings emerge. The fitted map has
$K{+}1=8$ factors: \Gfac{} (overall burden) and seven specific axes --- cognition,
immunometabolic, sleep, mania/activation, suicidality, developmental-risk and substance.
The fitted map carries $109$ indicators
($88$ in the Gaussian-copula continuous block, $21$ in the explicit non-Gaussian block).

\subsection{Mixed-likelihood observation model and the Gaussian copula}
\label{sub:mixed}

Coercing variables onto a common scale would manufacture covariance and bias the very
factors being estimated. Each indicator instead keeps the observation density appropriate
to its type, all sharing the linear predictor \eqref{eq:linpred}:
\begin{align}
  \text{continuous / $z$-score:}\quad & X_{ij}\sim \mathcal{N}(\eta_{ij},\ \sigma_j^2),
  \label{eq:fam-gauss}\\
  \text{skewed laboratory:}\quad & \log X_{ij}\sim \mathcal{N}(\eta_{ij},\ \sigma_j^2), \\
  \text{ordinal ($C_j$ levels):}\quad & \Pr(X_{ij}\le c)
  =\operatorname{logit}^{-1}(\tau_{jc}-\eta_{ij}),\ \ \tau_{j1}<\dots<\tau_{j,C_j-1},
  \label{eq:fam-ord}\\
  \text{binary:}\quad & X_{ij}\sim \mathrm{Bernoulli}\!\left(\operatorname{logit}^{-1}(\eta_{ij})\right),
  \label{eq:fam-bern}\\
  \text{count:}\quad & X_{ij}\sim \mathrm{NegBinom}\!\left(\mu_{ij}=e^{\eta_{ij}},\ r_j\right).
  \label{eq:fam-count}
\end{align}
The reported map replaces the raw Gaussian likelihood on the continuous block with a
semiparametric Gaussian-copula model: each continuous indicator is mapped to a latent
Gaussian score by its frozen empirical distribution function
$\hat F_j$, $z_{ij}=\Phi_{\mathcal N}^{-1}(\hat F_j(X_{ij}))$, and the factor model is fit
on $z_{ij}$. Because the rank transform is monotone it alters none of the identification
arguments below. Writing $F_j(\cdot)$ for indicator $j$'s density and
$O_i\subseteq\{1,\dots,J\}$ for its observed indicators, the observed-data log-likelihood
sums over observed cells only,
\begin{equation}
  \log p(X_{\mathrm{obs}}\mid\Theta)=\sum_{i=1}^{N}\sum_{j\in O_i}
  \log F_j\!\left(X_{ij}\mid\eta_{ij},\theta_j\right),
  \label{eq:obslik}
\end{equation}
which is \MEq{eq:obslik} of the main text: a missing cell contributes no term and no
completed $N\times J$ matrix is ever formed.

\subsection{Missing data is the marginal property, not imputation}
\label{sub:missing}

\begin{proposition}[Marginalization]\label{prop:marg}
For patient $i$ with observed continuous set $C_i\subseteq O_i$, the model-implied
distribution of the observed sub-vector is obtained by \emph{deleting} the unobserved
coordinates from $(\mu,\Sigma)$:
\begin{equation}
  X_{i,C_i}\sim\mathcal{N}\!\left(\mu_{i,C_i},\ \Sigma_{C_i}\right),\qquad
  \Sigma_{C_i}=\Lambda_{C_i}\Phi\Lambda_{C_i}^{\top}+\Psi_{C_i},
  \label{eq:margdist}
\end{equation}
with $\Lambda_{C_i}$ the observed-row loadings and
$\Psi_{C_i}=\operatorname{diag}(\sigma_j^2)_{j\in C_i}$.
\end{proposition}
\begin{proof}
A foundational property of the multivariate normal is that marginals are read off by
deleting coordinates: if $x\sim\mathcal N(\mu,\Sigma)$ then any sub-vector $x_{C}$ is
$\mathcal N(\mu_C,\Sigma_{CC})$. Apply this to $x_i\sim\mathcal N(\mu_i,\Sigma)$ with
$\Sigma=\Lambda\Phi\Lambda^\top+\Psi$ (Proposition~\ref{prop:sigma}); the submatrix
$\Sigma_{CC}$ retains exactly the observed rows and columns of $\Lambda$ and the observed
diagonal of $\Psi$, giving \eqref{eq:margdist}.
\end{proof}
\noindent
Hence ``integrating out'' the missing cells is not an extra step but the marginal property
itself; summing $\log\mathcal N(\cdot)$ over each patient's observed sub-vector is precisely
the full-information maximum-likelihood (FIML) objective for incomplete
multivariate-normal data. This is the formal justification of the no-imputation invariant
\eqref{eq:obslik}: any completed matrix would inject between-cell covariance that
\eqref{eq:localindep} would misattribute to a latent factor.

\subsection{Identifiability: rotation invariance and how the prior resolves it}
\label{sub:rotation}

\begin{proposition}[Rotation invariance of the likelihood]\label{prop:rotation}
For any invertible $M\in\mathbb{R}^{(K+1)\times(K+1)}$, the reparameterization
$\Lambda^{\star}=\Lambda M$, $\Phi^{\star}=M^{-1}\Phi M^{-\top}$, $f^{\star}=M^{-1}f$
leaves the implied covariance, and hence the Gaussian likelihood, unchanged:
\begin{equation}
  \Lambda^{\star}\Phi^{\star}\Lambda^{\star\top}+\Psi
  =\Lambda M\,(M^{-1}\Phi M^{-\top})\,M^{\top}\Lambda^{\top}+\Psi
  =\Lambda\Phi\Lambda^{\top}+\Psi=\Sigma .
  \label{eq:rotation}
\end{equation}
\end{proposition}
\begin{proof}
Immediate from $MM^{-1}=I$: the inner $M(M^{-1}\Phi M^{-\top})M^\top$ telescopes to $\Phi$.
Since $x_i$ is determined in distribution by $(\mu_i,\Sigma)$ and $\Sigma$ is invariant, so
is the likelihood.
\end{proof}
\noindent
A factor model is therefore not identified from the likelihood alone. Confirmatory factor
analysis breaks this by fixing loadings to exactly zero; the present model breaks it with
the prior. Each (indicator, factor) cell is assigned a role and prior
\begin{equation}
  \lambda_{jk}\mid\rho(j,k)\ \sim\
  \begin{cases}
    \mathcal{N}_{+}(0.70,\,0.25^2) & \textsc{primary (home cell)},\\[2pt]
    \text{regularized horseshoe } \eqref{eq:horseshoe} & \textsc{plausible cross},\\[2pt]
    \mathcal{N}(0,\,0.05^2) & \textsc{unlikely},\\[2pt]
    \delta_{0} & \textsc{forbidden},
  \end{cases}
  \label{eq:softprior}
\end{equation}
with $\mathcal{N}_{+}$ truncated to $(0,\infty)$ and $\delta_0$ fixing the loading at zero.
The likelihood is rotation-invariant but the prior is not: a rotation moves a loading the
prior wanted near zero away from zero, lowering $p(\Lambda)$, so the posterior
$p(\Lambda\mid X)\propto p(X\mid\Lambda)\,p(\Lambda)$ concentrates at the orientation where
small loadings align with the soft-zero priors and home loadings are positive and
substantial; the positive truncation removes the residual sign-flip. The plausible-cross
cells use a regularized (``Finnish'') horseshoe,
\begin{equation}
  \lambda_j=z_j\,\tau\,\tilde\eta_j,\qquad
  \tilde\eta_j=\sqrt{\dfrac{c^2\,\eta_j^2}{c^2+\tau^2\eta_j^2}},\qquad
  z_j\sim\mathcal N(0,1),\ \eta_j\sim\mathrm{HalfCauchy}(1),\ \tau\sim\mathrm{HalfNormal}(\tau_0),
  \label{eq:horseshoe}
\end{equation}
with $\tau_0=0.05$ and slab $c=0.30$: default-off (global $\tau$ shrinks the set toward
zero), evidence-on (heavy-tailed local $\eta_j$ lets a genuine cross-loading escape), and
magnitude-capped ($\tilde\eta_j\to c/\tau$, so an escaped loading stays small). This is the
honest treatment of weak signal in a factor with few indicators.

\begin{corollary}[Flat-prior confirmation]\label{cor:flat}
Under a flat loading prior (identification constraints only) the posterior mode equals the
maximum-likelihood (FIML) solution. A prior-free refit landing on the same loadings
therefore demonstrates that the structure is recovered from the data, not manufactured by
the priors; empirically this refit reproduced the map to Tucker congruence $\varphi=1.00$.
\end{corollary}

\subsection{The covariance identity and the correlated-\texorpdfstring{\Gfac}{G} estimand}
\label{sub:sigma}

\begin{proposition}[Bifactor covariance]\label{prop:sigma}
Under \eqref{eq:xstack} with $\operatorname{Var}(G_i)=1$,
$\operatorname{Cov}(D_i)=\Phi_D$, $\eED[\varepsilon_i\varepsilon_i^\top]=\Psi$ diagonal,
and the bifactor orthogonality $\eED[G_iD_i^\top]=\mathbf 0$ together with
$(G_i,D_i)\perp\varepsilon_i$,
\begin{equation}
  \Sigma=\operatorname{Cov}(x_i)
  =\underbrace{\lambda_G\lambda_G^{\top}}_{\text{general: rank one}}
  +\underbrace{\Lambda_D\Phi_D\Lambda_D^{\top}}_{\text{specific factors}}
  +\underbrace{\Psi}_{\text{private noise}}
  =\Lambda\,\Phi\,\Lambda^{\top}+\Psi .
  \label{eq:sigmaid}
\end{equation}
\end{proposition}
\begin{proof}
Expand $\Sigma=\eED[(x_i-\mu_i)(x_i-\mu_i)^\top]$ with
$x_i-\mu_i=\lambda_G G_i+\Lambda_D D_i+\varepsilon_i$. The cross-terms
$\lambda_G\,\eED[G_iD_i^\top]\,\Lambda_D^\top$ and
$\eED[(\lambda_GG_i+\Lambda_DD_i)\varepsilon_i^\top]$ vanish by
$\eED[G_iD_i^\top]=\mathbf 0$ and by independence of factors and noise respectively. The
survivors are $\lambda_G\lambda_G^\top\operatorname{Var}(G_i)$, $\Lambda_D\Phi_D\Lambda_D^\top$
and $\Psi$, giving \eqref{eq:sigmaid}.
\end{proof}
\noindent
The general factor leaves a rank-one footprint: every pair of items covaries through \Gfac{}
by exactly $\lambda_{jG}\lambda_{j'G}$. The correlated-\Gfac{} arm (\MEq{eq:corrg} of the
main text) unzips the zero off-diagonal block, freeing
$\phi_{Gd}=\operatorname{Corr}(G,D)$, whose immunometabolic entry near zero against far
larger cognition and sleep entries is the estimand the companion paper interprets. Here it
is used only as a worked example of the instrument.

%=============================================================================
\section{Projection: the Woodbury reduction and EAP scoring}
\label{sec:project}

\subsection{Woodbury reduction of the observed-cell covariance}
\label{sub:woodbury}

Evaluating \eqref{eq:margdist} naively costs $\mathcal{O}(|C_i|^3)$ per patient. The
Woodbury identity and matrix-determinant lemma reduce it to the $(K{+}1)$-dimensional
factor space.

\begin{lemma}[Woodbury reduction]\label{lem:woodbury}
With $\Psi_{C_i}$ diagonal and $\Lambda_{C_i}$ the observed-row loadings,
\begin{align}
  \Sigma_{C_i}^{-1} &= \Psi_{C_i}^{-1}-\Psi_{C_i}^{-1}\Lambda_{C_i}
      \bigl(\Phi^{-1}+\Lambda_{C_i}^{\top}\Psi_{C_i}^{-1}\Lambda_{C_i}\bigr)^{-1}
      \Lambda_{C_i}^{\top}\Psi_{C_i}^{-1}, \label{eq:woodbury}\\[2pt]
  \log\!\left|\Sigma_{C_i}\right| &=
      \log\!\left|\Phi^{-1}+\Lambda_{C_i}^{\top}\Psi_{C_i}^{-1}\Lambda_{C_i}\right|
      +\log|\Phi|+\!\!\sum_{j\in C_i}\!\log\sigma_j^2 .
      \label{eq:logdet}
\end{align}
\end{lemma}
\begin{proof}
Equation~\eqref{eq:woodbury} is the Woodbury matrix identity applied to
$\Sigma_{C_i}=\Psi_{C_i}+\Lambda_{C_i}\Phi\Lambda_{C_i}^{\top}$ with the low-rank term
$\Lambda_{C_i}\Phi\Lambda_{C_i}^{\top}$; \eqref{eq:logdet} is the matrix-determinant lemma
$\log|\Psi+U\Phi U^\top|=\log|\Phi^{-1}+U^\top\Psi^{-1}U|+\log|\Phi|+\log|\Psi|$ with
$U=\Lambda_{C_i}$ and $\Psi$ diagonal.
\end{proof}
\noindent
The inner solve is $\mathcal{O}((K{+}1)^3)$, independent of $|C_i|$; patients sharing a
missingness pattern share the inner $(K{+}1)\times(K{+}1)$ factorization, computed once per
pattern. This makes full-sample ($N=9{,}013$) projection tractable and is \MEq{eq:m-woodbury}
of the main text; the non-Gaussian block retains explicit per-patient latents.

\subsection{EAP factor scores and Bayes/FIML equivalence}
\label{sub:eap}

\begin{proposition}[EAP scores]\label{prop:eap}
For an all-continuous observed pattern, the posterior of $f_i$ given the observed cells is
Gaussian with mean and covariance
\begin{equation}
  \hat f_i=\eED[f_i\mid X_{i,C_i}]
  =\Phi\Lambda_{C_i}^{\top}\Sigma_{C_i}^{-1}\!\left(X_{i,C_i}-\mu_{i,C_i}\right),\qquad
  S_i=\Phi-\Phi\Lambda_{C_i}^{\top}\Sigma_{C_i}^{-1}\Lambda_{C_i}\Phi ,
  \label{eq:eap}
\end{equation}
which is \MEq{eq:m-scores} of the main text.
\end{proposition}
\begin{proof}
Under \eqref{eq:factors} and \eqref{eq:xstack} the joint of $(f_i,X_{i,C_i})$ is Gaussian
with $\operatorname{Cov}(f_i)=\Phi$,
$\operatorname{Cov}(X_{i,C_i})=\Sigma_{C_i}$ and cross-covariance
$\operatorname{Cov}(f_i,X_{i,C_i})=\Phi\Lambda_{C_i}^{\top}$. The Gaussian conditioning
formula $\eED[f\mid X]=\operatorname{Cov}(f,X)\operatorname{Cov}(X)^{-1}(X-\mu_X)$ and
$\operatorname{Cov}(f\mid X)=\operatorname{Cov}(f)-\operatorname{Cov}(f,X)
\operatorname{Cov}(X)^{-1}\operatorname{Cov}(X,f)$ give \eqref{eq:eap} directly.
\end{proof}

\begin{proposition}[Equivalent precision form]\label{prop:precform}
The posterior covariance \eqref{eq:eap} admits the precision form
\begin{equation}
  S_i^{-1}=\Phi^{-1}+\Lambda_{C_i}^{\top}\Psi_{C_i}^{-1}\Lambda_{C_i},
  \label{eq:precform}
\end{equation}
so the posterior precision is the prior precision plus a data term that is a sum of
per-item contributions.
\end{proposition}
\begin{proof}
Substitute the Woodbury form \eqref{eq:woodbury} into
$S_i=\Phi-\Phi\Lambda_{C_i}^{\top}\Sigma_{C_i}^{-1}\Lambda_{C_i}\Phi$ and simplify; the
low-rank correction telescopes to give $S_i=(\Phi^{-1}+\Lambda_{C_i}^{\top}\Psi_{C_i}^{-1}
\Lambda_{C_i})^{-1}$. Equivalently, \eqref{eq:precform} is the standard Gaussian
information form of the posterior of a linear model with prior precision $\Phi^{-1}$ and
per-observation precision $\Psi_{C_i}^{-1}$ propagated through $\Lambda_{C_i}$. Inverting
gives the covariance form \eqref{eq:eap}.
\end{proof}
\noindent
Equation~\eqref{eq:precform} is the pivot of the whole instrument: it exhibits the posterior
precision as \emph{additive over observed items}, which \S\ref{sec:info} turns into the
information accounting. When non-Gaussian cells are present the posterior is no longer exactly
Gaussian; $\hat f_i$ and $S_i$ are then the posterior mean and covariance obtained by MCMC or,
for the calibration study, by the Fisher-scoring Laplace approximation of \S\ref{sec:calib},
which reduces to \eqref{eq:eap} exactly on all-Gaussian patterns.

\begin{proposition}[Bayes/FIML equivalence]\label{prop:fiml}
The marginalized Bayesian observed-data model and FIML optimize the same observed-data
likelihood \eqref{eq:obslik} and differ only by the priors \eqref{eq:softprior}. In
particular (Corollary~\ref{cor:flat}) the flat-prior MAP equals the FIML estimate, so a
standalone FIML estimator adds no independent evidence and confirmation is done in-engine.
\end{proposition}
\begin{proof}
By Proposition~\ref{prop:marg} the per-patient marginal density is the incomplete-data
multivariate-normal density on $C_i$, whose log summed over patients is the FIML objective.
The posterior is this likelihood times the prior; with a flat prior the log-posterior equals
the log-likelihood up to a constant, so its mode is the FIML maximizer.
\end{proof}
\noindent
Finally, the structure of \eqref{eq:eap} makes the instrument's honesty on unmeasured axes
mechanical: on an axis with no observed indicator the corresponding rows of $\Lambda_{C_i}$
are empty, the update term in \eqref{eq:precform} vanishes on that coordinate, and the
posterior returns the prior $\mathcal N(0,\Phi)$ --- a wide interval, never a false-precise
point. The per-patient posterior covariance $S_i$ is exported and propagated into every
downstream analysis; a patient with one observed indicator on a factor is not treated as
equally characterized as one with six.

%=============================================================================
\section{The loading as a measurement slope}
\label{sec:slope}

The main text (\MSec{sec:m-slope}) separates two objects that \eqref{eq:linpred} places on
the same line: the per-patient latent \emph{position} $f_i$, and the fixed \emph{slope} at
which a measurement changes as that position moves. This section derives the slope on both
the latent and the response scale.

\begin{proposition}[Latent-scale slope]\label{prop:latslope}
On the linear-predictor scale the sensitivity of indicator $j$ to factor $k$ is exact,
constant across patients, and equal to the loading:
\begin{equation}
  \frac{\partial\eta_{ij}}{\partial f_{ik}}=\lambda_{jk}.
  \label{eq:latslope}
\end{equation}
Hence the loading matrix $\Lambda$ is the Jacobian $\partial\eta_i/\partial f_i$ of the
measurement with respect to the latent coordinate.
\end{proposition}
\begin{proof}
Differentiate \eqref{eq:linpred} with respect to $f_{ik}$; the intercept, the covariate
term $\beta_j^\top c_i$ and the other factor terms do not depend on $f_{ik}$, leaving
$\lambda_{jk}$.
\end{proof}

\begin{proposition}[Response-scale slope]\label{prop:respslope}
For a single-index likelihood with conditional mean
$\eED[X_{ij}\mid f_i]=h_j(\eta_{ij})$ (link mean function $h_j$), the marginal effect of a
factor on the indicator and its population average are
\begin{equation}
  \frac{\partial\,\eED[X_{ij}\mid f_i]}{\partial f_{ik}}=h_j'(\eta_{ij})\,\lambda_{jk},
  \qquad
  s_{jk}\equiv\eED\!\left[\frac{\partial\,\eED[X_{ij}\mid f_i]}{\partial f_{ik}}\right]
  =\kappa_j\,\lambda_{jk},\qquad
  \kappa_j=\eED\!\left[h_j'(\eta_{ij})\right]\ge 0 ,
  \label{eq:respslope}
\end{equation}
which is \MEq{eq:m-slope} of the main text.
\end{proposition}
\begin{proof}
The chain rule applied to $\eED[X_{ij}\mid f_i]=h_j(\eta_{ij})$ with
$\partial\eta_{ij}/\partial f_{ik}=\lambda_{jk}$ (Proposition~\ref{prop:latslope}) gives the
first identity. Taking the expectation over the population distribution of $\eta_{ij}$ and
factoring out the patient-independent $\lambda_{jk}$ defines
$\kappa_j=\eED[h_j'(\eta_{ij})]$, which is non-negative because each family's mean function
$h_j$ is non-decreasing ($h_j'=\sigma^{-2}$-scaled identity for the Gaussian,
$p(1-p)$ for the logit-Bernoulli, $e^{\eta}$ for the log-count).
\end{proof}
\noindent
The loading is therefore \emph{how strongly an item reflects a factor}. It is not how much
that item localizes a patient --- that is the Fisher information of \S\ref{sec:info}, and the
two coincide only for the Gaussian family, where $h_j'$ is constant. The response-scale
average $\kappa_j$ is what makes a unit-loading rare binary flag a weak \emph{measurement}
despite a large $\lambda$: $\kappa_j=\eED[p_{ij}(1-p_{ij})]$ collapses as the endorsement
probability leaves $\tfrac12$.

%=============================================================================
\section{Information accounting: what localizes a patient}
\label{sec:info}

\subsection{Additivity of posterior precision}
\label{sub:additive}

\begin{proposition}[Additive precision]\label{prop:additive}
Under local independence \eqref{eq:localindep} the posterior precision of the coordinate is
the prior precision plus a sum of per-item Fisher-information contributions,
\begin{equation}
  S_i^{-1}=\Phi^{-1}+\sum_{j\in O_i}\info_j(f_i),
  \qquad
  \info_j(f_i)=-\,\eED_{X_{ij}\mid f_i}\!\left[\nabla^2_{f_i}\log F_j(X_{ij}\mid\eta_{ij})\right],
  \label{eq:precision}
\end{equation}
which is \MEq{eq:m-precision} of the main text. Each observed cell adds a positive
semidefinite term; nothing an item measures is ever subtracted.
\end{proposition}
\begin{proof}
The log-posterior of $f_i$ given the observed cells is, up to a constant,
$\log p(f_i\mid X_{i,O_i})=\log\mathcal N(f_i;\mathbf 0,\Phi)
+\sum_{j\in O_i}\log F_j(X_{ij}\mid\eta_{ij})$, the sum over items following from local
independence \eqref{eq:localindep}. The negative Hessian in $f_i$ of the prior term is
$\Phi^{-1}$; the negative expected Hessian of each item term is by definition the item's
Fisher information $\info_j(f_i)$ about the factors. Summing gives \eqref{eq:precision}; each
$\info_j$ is positive semidefinite as an expected outer product of score vectors (below), so
observing an item can only increase precision. For the Gaussian family the expectation is
exact and \eqref{eq:precision} coincides with \eqref{eq:precform}; for the other families it
is the expected-information (Fisher-scoring) form used in \S\ref{sec:calib}.
\end{proof}

\subsection{Per-family unit information}
\label{sub:family}

Because $\eta_{ij}$ is a single index in $f_i$ with $\nabla_{f_i}\eta_{ij}=\lambda_j$
(Proposition~\ref{prop:latslope}), the information that item $j$ carries about the factors
factorizes through the scalar index:
\begin{equation}
  \info_j(f_i)=\unitinfo_j(\eta_{ij})\,\lambda_j\lambda_j^{\top},
  \qquad
  \unitinfo_j(\eta)=-\,\eED_{X\mid\eta}\!\left[\frac{\partial^2}{\partial\eta^2}\log F_j(X\mid\eta)\right],
  \label{eq:unitinfo}
\end{equation}
so the per-factor information is $\info_{jk}(f_i)=\lambda_{jk}^{2}\,\unitinfo_j(\eta_{ij})$ ---
\MEq{eq:m-fisher} of the main text --- with $\unitinfo_j$ the one-dimensional Fisher
information of the item's own likelihood family. We derive $\unitinfo_j$ for each family.

\begin{proposition}[Family unit informations]\label{prop:families}
For the five families of \eqref{eq:fam-gauss}--\eqref{eq:fam-count},
\begin{align}
  \text{Gaussian:}\quad & \unitinfo_j(\eta)=\frac{1}{\sigma_j^2}, \label{eq:info-gauss}\\
  \text{Bernoulli (logit):}\quad & \unitinfo_j(\eta)=p(1-p),\quad p=\operatorname{logit}^{-1}(\eta), \label{eq:info-bern}\\
  \text{Negative-binomial (log link):}\quad & \unitinfo_j(\eta)=\frac{\mu\,r_j}{\mu+r_j},\quad \mu=e^{\eta}, \label{eq:info-nb}\\
  \text{Graded-response ordinal:}\quad & \unitinfo_j(\eta)=\sum_{c=1}^{C_j}\frac{\bigl(\partial_\eta P_{jc}(\eta)\bigr)^2}{P_{jc}(\eta)}
    =\sum_{c=1}^{C_j-1}\frac{\bigl(P_{jc}^{\star}(1-P_{jc}^{\star})-P_{j,c+1}^{\star}(1-P_{j,c+1}^{\star})\bigr)^2}{P_{jc}(\eta)}, \label{eq:info-ord}
\end{align}
where $P_{jc}(\eta)=P_{jc}^{\star}-P_{j,c+1}^{\star}$ is the probability of category $c$ and
$P_{jc}^{\star}=\operatorname{logit}^{-1}(\eta-\tau_{jc})$ is the cumulative probability of
exceeding threshold $\tau_{jc}$ (with $P_{j0}^{\star}=1$, $P_{j,C_j}^{\star}=0$).
\end{proposition}
\begin{proof}
\emph{Gaussian.} $\log F_j=-\tfrac12\log(2\pi\sigma_j^2)-\tfrac{(X-\eta)^2}{2\sigma_j^2}$, so
$\partial_\eta^2\log F_j=-1/\sigma_j^2$, a constant; its negative expectation is
$1/\sigma_j^2$, independent of $\eta$.

\emph{Bernoulli.} With $p=\operatorname{logit}^{-1}(\eta)$,
$\log F_j=X\log p+(1-X)\log(1-p)$ and $\partial_\eta p=p(1-p)$. The score is
$\partial_\eta\log F_j=X-p$, and
$\unitinfo_j=\operatorname{Var}(X-p\mid\eta)=\operatorname{Var}(X)=p(1-p)$.

\emph{Negative-binomial.} For $X\sim\mathrm{NegBinom}$ with mean $\mu=e^\eta$ and dispersion
$r_j$, $\operatorname{Var}(X\mid\eta)=\mu+\mu^2/r_j$. With the log link
$\partial_\eta\mu=\mu$, the generalized-linear-model information is
$\unitinfo_j=(\partial_\eta\mu)^2/\operatorname{Var}(X)=\mu^2/(\mu+\mu^2/r_j)=\mu r_j/(\mu+r_j)$.

\emph{Graded-response ordinal.} For a categorical likelihood
$\log F_j=\sum_c \mathbf 1\{X=c\}\log P_{jc}(\eta)$, the score is
$\partial_\eta\log F_j=\sum_c\mathbf 1\{X=c\}\,\partial_\eta P_{jc}/P_{jc}$ and the expected
squared score is $\unitinfo_j=\sum_c(\partial_\eta P_{jc})^2/P_{jc}$, the standard
graded-response information \citep{samejima1969}. Writing category probabilities as
differences of the cumulative logits $P_{jc}^\star$ and using
$\partial_\eta P_{jc}^\star=P_{jc}^\star(1-P_{jc}^\star)$ gives the closed form on the right
of \eqref{eq:info-ord}.
\end{proof}

\subsection{Why loading is not information}
\label{sub:notloading}

Equation~\eqref{eq:unitinfo} is why localization and reflection part company. The loading
$\lambda_{jk}$ enters $\info_{jk}$ squared, but it multiplies the unit information
$\unitinfo_j(\eta_{ij})$, which for every non-Gaussian family depends on where the patient
sits. The Bernoulli case \eqref{eq:info-bern} is the sharpest: a binary flag endorsed by a
fraction $p$ of patients contributes $\unitinfo_j=p(1-p)$, which is $0.09$ at $p=0.1$ and
collapses to $0$ as $p\to0$, \emph{regardless of $\lambda$}. A strongly-loading, rarely-endorsed
symptom therefore localizes almost no one. Only for the Gaussian family
\eqref{eq:info-gauss}, where $\unitinfo_j=1/\sigma_j^2$ is constant, does the ranking by
information coincide with the ranking by (scaled) loading.

\begin{corollary}[Reliability and the target-driven item count]\label{cor:reliability}
On a single axis $k$, accumulating \eqref{eq:precision} gives a scalar precision
$\pi_k=1+\sum_{j}\info_{jk}$ (prior precision $1$ for the standardized factor) and a
reliability
\begin{equation}
  \rho_k=1-\frac{1}{\pi_k}=1-\frac{1}{\,1+\sum_j \lambda_{jk}^2\,\unitinfo_j\,}.
  \label{eq:reliability}
\end{equation}
To reach a posterior standard deviation $\pi_k^{-1/2}\le\varsigma$ requires accumulated
information $\sum_j\info_{jk}\ge \varsigma^{-2}-1$; at $\varsigma=0.5$ this is
$\sum_j\info_{jk}\ge 3$. For a standardized continuous indicator of common loading
$\lambda$ the residual variance is $1-\lambda^2$, so its Gaussian unit information is
$\unitinfo=1/(1-\lambda^2)$ and it contributes $\lambda^2/(1-\lambda^2)$ per item; the
required count is therefore
\begin{equation}
  n \ \ge\ \frac{\varsigma^{-2}-1}{\lambda^2/(1-\lambda^2)}
  \ =\ 3\,\frac{1-\lambda^2}{\lambda^2}\quad(\varsigma=0.5),
  \label{eq:itemcount}
\end{equation}
i.e.\ $3.1$ items at $\lambda=0.7$, $9.0$ at $\lambda=0.5$ and $30.3$ at $\lambda=0.3$ ---
consistent with the ``roughly $4$, $9$ and $31$ items'' the full mixed-bank simulation
reports in \MSec{sec:res-info}.
\end{corollary}
\begin{proof}
Immediate from $\pi_k=1/\operatorname{Var}(f_{ik}\mid X)$ and \eqref{eq:reliability}: the
item count solves $\lambda^2\unitinfo\cdot n=\varsigma^{-2}-1$ with the standardized-indicator
unit information $\unitinfo=1/(1-\lambda^2)$, giving \eqref{eq:itemcount}.
\end{proof}

%=============================================================================
\section{Adaptive sampling and value of information}
\label{sec:voi}

\subsection{The fixed order as an upper bound on adaptive efficiency}
\label{sub:upperbound}

Ordering the candidate items for an axis by their information \eqref{eq:unitinfo} and
accumulating \eqref{eq:precision} gives the reliability-versus-items curve
\eqref{eq:reliability}. Evaluated at the population mean $f_i=\mathbf 0$ this is the
efficiency of a fixed, non-adaptive most-informative-first battery. A computerized adaptive
test (CAT) instead re-evaluates $\unitinfo_j(\eta_{ij})$ at each patient's provisional
estimate $\hat\theta$ and selects the currently most informative item.

\begin{proposition}[Population-mean order is an upper bound]\label{prop:upperbound}
Fix an axis and let $\bar{\info}_j=\unitinfo_j(\bar\eta_j)$ be the unit information evaluated
at the population-mean linear predictor $\bar\eta_j=\alpha_j$. The fixed schedule that
administers items in decreasing $\lambda_j^2\bar\info_j$ maximizes accumulated information at
every step among schedules that do not use the patient's responses. Any adaptive schedule
attains, in expectation over patients, no more accumulated information per item at the
population mean than this fixed order.
\end{proposition}
\begin{proof}
At the population mean the per-item increments $\lambda_j^2\bar\info_j$ are fixed constants;
by the rearrangement inequality, administering them in decreasing order maximizes the
partial sums $\sum_{j\le n}$ at every $n$, so the fixed most-informative-first order is
optimal among response-independent schedules. An adaptive schedule can exceed a fixed one
only by exploiting the dependence of $\unitinfo_j(\eta_{ij})$ on the unknown $f_i$; at the
population mean, and for families whose $\unitinfo_j$ does not depend on $\eta$ (the Gaussian
family, \eqref{eq:info-gauss}), that dependence is absent and the two schedules coincide.
\end{proof}
\noindent
This is why the CAT simulation of \MSec{sec:res-adaptive} finds the bound \emph{tight}: on
every axis whose leading items are Gaussian the adaptive and fixed curves coincide to within
$2\times10^{-4}$ in reliability, because $\unitinfo_j$ is constant and there is nothing to
adapt to; only Bernoulli-topped axes (whose $\unitinfo_j=p(1-p)$ varies with $\theta$) show a
small adaptive signature, at most $0.012$ in reliability. The bank-limited axes are a
separate phenomenon: mania/activation and substance plateau at reliability $0.41$ and $0.43$
because $\sum_j\lambda_{jk}^2\unitinfo_j$ saturates over a $2$- and $4$-item bank
respectively --- a ceiling of \eqref{eq:reliability} set by instrument content, not by any
schedule.

\subsection{Value of information: the matrix-determinant lemma}
\label{sub:voi}

Across axes the same information becomes a design criterion through the entropy of the
posterior coordinate.

\begin{proposition}[Closed-form entropy reduction]\label{prop:voi}
The differential entropy of the Gaussian posterior is
$H(f_i)=\tfrac12\log\det(2\pi e\,S_i)$. Adding a single not-yet-measured indicator
$j^\star$ with loading $\lambda_{j^\star}$ and unit information
$\unitinfo_{j^\star}=\unitinfo_{j^\star}(\hat\eta_{ij^\star})$ reduces the entropy by
\begin{equation}
  H(f_i)-H(f_i\mid j^\star)
  =\tfrac12\log\!\left(1+\unitinfo_{j^\star}\,\lambda_{j^\star}^{\top} S_i\,\lambda_{j^\star}\right),
  \label{eq:voi}
\end{equation}
a closed form that depends on what has already been measured through the current posterior
covariance $S_i$. Equation~\eqref{eq:voi} is \MEq{eq:m-voi} of the main text.
\end{proposition}
\begin{proof}
By Proposition~\ref{prop:additive} the updated precision is
$S_i^{-1}\!\mid_{+j^\star}=S_i^{-1}+\unitinfo_{j^\star}\lambda_{j^\star}\lambda_{j^\star}^{\top}$,
a rank-one update. The entropy change is
$H(f_i)-H(f_i\mid j^\star)=\tfrac12\log\det(S_i^{-1}\!\mid_{+j^\star})-\tfrac12\log\det(S_i^{-1})
=\tfrac12\log\frac{\det(S_i^{-1}+\unitinfo_{j^\star}\lambda_{j^\star}\lambda_{j^\star}^{\top})}{\det S_i^{-1}}$.
The matrix-determinant lemma $\det(A+uv^\top)=\det(A)(1+v^\top A^{-1}u)$ with $A=S_i^{-1}$,
$u=\unitinfo_{j^\star}\lambda_{j^\star}$, $v=\lambda_{j^\star}$ gives the ratio
$1+\unitinfo_{j^\star}\lambda_{j^\star}^{\top}S_i\lambda_{j^\star}$, and the logarithm halves
to \eqref{eq:voi}.
\end{proof}
\noindent
Two properties make \eqref{eq:voi} the right design criterion. It is \emph{context-dependent}:
the value of a candidate item scales with $\lambda_{j^\star}^{\top}S_i\lambda_{j^\star}$, the
current posterior variance along the item's loading direction, so an item is worth most on an
axis that is still uncertain and its value falls as that axis fills in. And it is
\emph{diminishing}: because $S_i$ shrinks with each added item, the marginal entropy reduction
of the next item decreases, so a greedy rule spreads measurement across axes rather than
piling onto one. Cycling the greedy selection across the eight axes assembles the shared
$27$-item battery at mean reliability $0.70$ of \MSec{sec:res-voi}; applied within a cohort,
the same $S_i$-weighted criterion localizes where an axis is under-measured (e.g.\ the
schizophrenia immunometabolic and sleep deficits, $20.7$ vs $34.9$ and $3.9$ vs $8.4$ items).

\begin{corollary}[Greedy value is a lower bound on the jointly-optimal design]\label{cor:greedy}
The entropy reduction $f\mapsto\log\det(\Phi^{-1}+\sum_{j\in S}\info_j)$ as a function of the
selected set $S$ is monotone and submodular, so the greedy battery attains at least a
$(1-1/e)$ fraction of the entropy reduction of the jointly-optimal same-size battery.
\end{corollary}
\begin{proof}
Monotonicity holds because each $\info_j\succeq0$ (Proposition~\ref{prop:additive}) can only
increase the log-determinant. Submodularity of $S\mapsto\log\det(A+\sum_{j\in S}\info_j)$ for
positive semidefinite $\info_j$ is the standard concavity-of-log-det result: the marginal gain
of adding $j^\star$, given by \eqref{eq:voi}, is non-increasing as the base set grows because
$S_i$ (the inverse of the accumulated precision) is decreasing in the Loewner order. The
$(1-1/e)$ guarantee is the Nemhauser–Wolsey–Fisher bound for greedy maximization of a
monotone submodular function. This bounds the price of greediness noted as a limitation in
the main text.
\end{proof}

%=============================================================================
\section{Archetypes as a low-dimensional summary and its reconstruction fidelity}
\label{sec:arch}

\subsection{The archetypal-analysis program}
\label{sub:aa}

The cloud of posterior coordinates $\{x_i\}_{i=1}^N$, $x_i=\hat f_i\in\mathbb R^{K+1}$, is
summarized by archetypal analysis \citep{cutler1994archetypal}: each patient is written as a
convex blend of $A$ extreme phenotypes that are themselves convex combinations of patients,
\begin{equation}
  \min_{W,B}\ \sum_{i=1}^{N}\Big\|x_i-\sum_{a=1}^{A} w_{ia}\,z_a\Big\|^{2}
  \quad\text{s.t.}\quad z_a=\sum_{j=1}^{N} b_{ja}\,x_j,\ \
  w_{ia},b_{ja}\ge0,\ \ \textstyle\sum_a w_{ia}=\sum_j b_{ja}=1 ,
  \label{eq:aa}
\end{equation}
i.e.\ $X\approx WZ$ with archetypes $Z=BX$ and both $W$ (patient$\times$archetype) and $B$
(patient$\times$archetype) row-stochastic. This is \MEq{eq:m-archetype} of the main text. The
constraints place each archetype $z_a$ inside the convex hull of the data and each patient's
reconstruction $\hat x_i=\sum_a w_{ia}z_a$ inside the convex hull of the archetypes; the
simplex weight $w_i\in\Delta^{A-1}$ is the patient's soft membership. The number of corners is
fixed by stability: $A=5$ is the largest with cross-seed Tucker congruence $\ge0.8$, with a
cliff at $A=6$, and $w_i$ is propagated across posterior draws.

\subsection{Reconstruction fidelity as a variance decomposition}
\label{sub:recon}

Because \eqref{eq:aa} is a rank-$A$ convex approximation it discards information, and unlike
an unconstrained projection its residual is not orthogonal to the fitted part. We quantify the
loss by a coefficient of determination.

\begin{definition}[Reconstruction $R^2$]\label{def:r2}
With reconstruction $\hat x_i=\sum_a w_{ia}z_a$, the per-axis and pooled reconstruction
$R^2$ are
\begin{equation}
  R^2_k=1-\frac{\sum_i (x_{ik}-\hat x_{ik})^2}{\sum_i (x_{ik}-\bar x_k)^2},
  \qquad
  R^2_{\mathrm{pool}}=1-\frac{\sum_i\|x_i-\hat x_i\|^2}{\sum_i\|x_i-\bar x\|^2}
  =\frac{\sum_k \mathrm{SS}^{\mathrm{tot}}_k R^2_k}{\sum_k \mathrm{SS}^{\mathrm{tot}}_k},
  \label{eq:r2}
\end{equation}
the pooled value being the variance-weighted mean of the per-axis values with weights
$\mathrm{SS}^{\mathrm{tot}}_k=\sum_i (x_{ik}-\bar x_k)^2$.
\end{definition}
\noindent
Evaluated on all $N=9{,}013$ patients the five-corner blend attains
$R^2_{\mathrm{pool}}=0.590$: the convex summary retains $59\%$ of the coordinate variance. It is
tempting to read this against the unconstrained linear optimum at the same dimension --- a
five-component principal-component projection attains $R^2=0.796$, a dimension-matched
four-component projection $0.680$ --- and conclude that convex interpretability costs
$20.6$ percentage points of fidelity. That reading is wrong, and Proposition~\ref{prop:kd} below
explains why: on this map the coordinate variance is nearly isotropic, so reconstruction $R^2$ is
an uninformative test of any five-dimensional summary. A \emph{random} five-dimensional subspace
retains an expected $k/d=5/8=0.625$ of the variance regardless of the covariance spectrum --- more
than the archetype simplex --- and a $500$-seed simulation confirms it ($0.624\pm0.042$, placing
the archetype value at the $21$st percentile of random projections). The simplex does not merely
trail PCA; it trails a coin-flip subspace, because variance capture is not what it optimizes.
What it does provide, and what neither PCA nor a random subspace can, is convex non-negative
sum-to-one membership in $A$ \emph{named} corners: a principal component is a signed axis-mixing
contrast that no patient can equal, whereas each archetype $z_a$ is a near-real extreme profile a
patient can be a stated fraction of. Reconstruction $R^2$ is therefore reported here as an honesty
check on how much of the continuum the corners span --- not as a competitive metric --- and the
five-centroid $k$-means hard partition ($R^2=0.325$), the discrete alternative appropriate only if
the cloud had clusters, is included to show that the least defensible summary is the categorical
one. Per axis, $R^2_k$ is high where the corners separate (immunometabolic $0.89$, overall
severity $0.86$, sleep $0.76$, developmental risk $0.75$) and near-flat on the two bank-limited
axes (cognition $0.06$, substance $\approx0$), the same axes flagged by \eqref{eq:reliability} ---
so the per-axis profile is a diagnostic of where the instrument is weak, not a scorecard.

\begin{proposition}[A random $k$-subspace captures $k/d$ of the variance]\label{prop:kd}
Let $x_1,\dots,x_N\in\mathbb R^d$ have mean zero and total variance
$T=\tfrac1N\sum_i\|x_i\|^2$. Let $P$ be the orthogonal projector onto a $k$-dimensional subspace
drawn uniformly (Haar) from the Grassmannian $\mathrm{Gr}(k,d)$. Then the expected variance
captured by projecting the data onto that subspace is
\begin{equation}
  \mathbb E\!\left[\tfrac1N\textstyle\sum_i\|Px_i\|^2\right]=\frac{k}{d}\,T,
  \qquad\text{i.e.}\qquad \mathbb E[R^2_{\text{rand-}k}]=\frac kd,
  \label{eq:kd}
\end{equation}
\emph{independent of the eigenvalue spectrum} of the data covariance.
\end{proposition}
\begin{proof}
Write the covariance $\Sigma=\tfrac1N\sum_i x_ix_i^\top$ with eigenvalues $\sigma_1^2\ge\dots\ge
\sigma_d^2$ summing to $T$. The captured variance is
$\tfrac1N\sum_i x_i^\top P x_i=\operatorname{tr}(P\Sigma)$. For a Haar-random rank-$k$ projector,
$\mathbb E[P]=(k/d)I_d$ by rotational invariance (the only symmetric matrix commuting with all
rotations is a scalar multiple of the identity, and $\operatorname{tr}\mathbb E[P]=k$ fixes the
scalar). Hence $\mathbb E[\operatorname{tr}(P\Sigma)]=\operatorname{tr}((k/d)\Sigma)=(k/d)T$, and
dividing by $T$ gives \eqref{eq:kd}. The result uses only $\mathbb E[P]=(k/d)I$; the spectrum
$\{\sigma_j^2\}$ cancels, which is why an anisotropic (e.g.\ single-dominant-component) cloud and a
near-isotropic one both give $k/d$ in expectation, while their PCA-optimal $R^2$ differ sharply.
\end{proof}
\noindent
For this map $k/d=5/8=0.625$; the simulated mean $0.624$ matches to Monte-Carlo error, and the
gap between that and the PCA-5 value $0.796$ is a direct measure of how far the coordinate cloud
departs from isotropy --- here, modestly (the leading component carries $0.269$ of the variance,
against $0.125$ for an exactly isotropic cloud).

\begin{proposition}[Fidelity does not degrade for ambiguous patients]\label{prop:fog}
Let $e_i=\|x_i-\hat x_i\|$ be a patient's reconstruction error and
$H_i=-\sum_a w_{ia}\log w_{ia}$ the entropy of their archetype weights. Then
$\operatorname{Corr}(e_i,H_i)<0$: interior ``fog'' patients reconstruct \emph{better}, not
worse.
\end{proposition}
\begin{proof}[Argument]
By construction each $\hat x_i$ lies in the convex hull of the archetypes. A high-entropy
patient has weights spread across corners and therefore sits near the centroid of the hull,
where the convex approximation is tight; a low-entropy (near-corner) patient sits near a
single archetype $z_a$, but the true $x_i$ near a corner frequently lies \emph{outside} the
hull spanned by the finite archetype set (the archetypes are interior convex combinations of
data, not the extreme points themselves), so the clipped reconstruction incurs a larger
residual. Empirically $\operatorname{Corr}(e_i,H_i)=-0.13$ ($p<10^{-30}$), consistent with
this geometry. The negative sign is the useful conclusion: the summary is not purchased by
sacrificing the ambiguous patients it is meant to describe.
\end{proof}

%=============================================================================
\section{Calibration protocol}
\label{sec:calib}

The instrument reports posterior intervals; \MSec{sec:m-calib} tests whether those intervals
are honest. Because the truth is unknown for real patients, coverage is tested by simulation
under the model's own generative law with real missingness.

\subsection{Synthetic-truth generative design}
\label{sub:gen-calib}

We draw $M=10{,}000$ synthetic patients (seed $20260704$) with known latent coordinate
$\theta_m\sim\mathcal N(\mathbf 0,\Phi)$ under the frozen model's factor prior. For each we
generate item responses from the fitted families
\eqref{eq:fam-gauss}--\eqref{eq:fam-count} at $\eta_{mj}=\alpha_j+\lambda_j^{\top}\theta_m$,
and impose a real missingness pattern by bootstrapping the observed/missing masks of the
$9{,}013$ real patients --- only the binary observed-or-missing indicator is used, never any
clinical value. This preserves the empirical joint distribution of \emph{which} axes are
co-measured, which is what stresses the projection.

\subsection{Fisher-scoring EAP and its exactness on Gaussian patterns}
\label{sub:fisher-eap}

Each synthetic patient is projected by the same machinery as a real one. For non-Gaussian
patterns the posterior mode and curvature are found by Fisher scoring: starting from
$\hat\theta^{(0)}=\mathbf 0$, iterate
\begin{equation}
  \hat\theta^{(t+1)}=\hat\theta^{(t)}
  +\Big(\Phi^{-1}+\textstyle\sum_{j\in O_m}\info_j(\hat\theta^{(t)})\Big)^{-1}
  \Big(\!-\Phi^{-1}\hat\theta^{(t)}+\textstyle\sum_{j\in O_m} u_{mj}(\hat\theta^{(t)})\Big),
  \label{eq:fisherscore}
\end{equation}
where $u_{mj}=\lambda_j\,\partial_\eta\log F_j(X_{mj}\mid\eta_{mj})$ is the item score and
$\info_j$ its expected information \eqref{eq:unitinfo}; the posterior covariance is the inverse
information $S_m=(\Phi^{-1}+\sum_{j}\info_j(\hat\theta))^{-1}$ at convergence, the Laplace
approximation to the posterior.

\begin{proposition}[Exactness on Gaussian patterns]\label{prop:exact}
If all observed items of patient $m$ are Gaussian \eqref{eq:fam-gauss}, the Fisher-scoring
iteration \eqref{eq:fisherscore} converges in one step from any start to the closed-form EAP
\eqref{eq:eap}, and $S_m$ equals the closed-form posterior covariance exactly.
\end{proposition}
\begin{proof}
For a Gaussian item $\info_j=\lambda_j\lambda_j^\top/\sigma_j^2$ is independent of $\theta$
and the score is affine, $u_{mj}(\theta)=\lambda_j(X_{mj}-\alpha_j-\lambda_j^\top\theta)/\sigma_j^2$.
The bracketed Newton system \eqref{eq:fisherscore} is then exactly the normal equations of the
Gaussian posterior, whose solution is \eqref{eq:precform}--\eqref{eq:eap}; a Newton step on a
quadratic objective reaches the optimum in one iteration regardless of the starting point.
\end{proof}
\noindent
We verified this numerically: on all-Gaussian patterns the Fisher-scoring coordinate matches
the closed form \eqref{eq:eap} to $3\times10^{-15}$, so the calibration study exercises the
\emph{same} estimator the instrument uses, extended to the non-Gaussian families.

\subsection{Coverage statistics and the prior-in-prior-out check}
\label{sub:coverage}

For nominal levels $q\in\{0.50,0.80,0.90,0.95\}$ with standard-normal quantile $z_q$, we
report three statistics: the per-axis marginal coverage
\begin{equation}
  \widehat{\mathrm{cov}}_k(q)=\frac{1}{M}\sum_{m=1}^{M}
  \mathbf 1\!\left\{\,|\theta_{mk}-\hat x_{mk}|\le z_{(1+q)/2}\,[S_m]_{kk}^{1/2}\right\},
  \label{eq:covax}
\end{equation}
the joint coverage of the $8$-dimensional credible ellipsoid via the Mahalanobis distance
against a $\chi^2_{K+1}$ quantile,
$\widehat{\mathrm{cov}}_{\mathrm{joint}}(q)=\tfrac1M\sum_m\mathbf 1\{(\theta_m-\hat x_m)^\top
S_m^{-1}(\theta_m-\hat x_m)\le \chi^2_{K+1,q}\}$, and the $95\%$ coverage stratified by the
number of observed indicators on the axis. The Monte-Carlo standard error of a coverage
estimate at level $q$ is $\sqrt{q(1-q)/M}$; at $q=0.95$, $M=10^4$ this is $0.0022$.

\begin{proposition}[Prior in, prior out]\label{prop:priorout}
On an axis $k$ with no observed indicator, the projection returns the prior on that
coordinate --- $\hat x_{mk}=0$ and $[S_m]_{kk}$ equal to the prior variance --- so the coverage
\eqref{eq:covax} equals the nominal level exactly, up to Monte-Carlo error.
\end{proposition}
\begin{proof}
With no observed item loading on axis $k$, the corresponding rows of $\Lambda_{O_m}$ are zero,
so the data term $\sum_j\info_j$ in \eqref{eq:precision} has no entry on coordinate $k$ and the
posterior marginal of $\theta_{mk}$ equals its prior $\mathcal N(0,\Phi_{kk})$. Since the
synthetic truth is drawn from that same prior \eqref{eq:factors}, the interval
$0\pm z_{(1+q)/2}\Phi_{kk}^{1/2}$ covers $\theta_{mk}$ with probability exactly $q$.
\end{proof}
\noindent
This makes the zero-item stratum a built-in correctness check rather than a free parameter:
the instrument \emph{must} recover nominal coverage there, and it does. Across all strata the
empirical coverage tracks nominal --- $0.501$, $0.801$, $0.900$, $0.949$ against
$0.50/0.80/0.90/0.95$, joint ellipsoid $0.947$ at $95\%$ --- and is flat from the
prior-dominated stratum to the richly-observed one ($0.947$–$0.952$ across item-count bins),
the result reported in \MSec{sec:res-calib}.

\paragraph{What the study does and does not establish.} Coverage here demonstrates that the
projection machinery is \emph{internally honest}: its $95\%$ intervals cover the truth $95\%$
of the time when the data are generated by the fitted model under real missingness. It does
not establish that the fitted model is the true data-generating process --- no coverage study
can, since it conditions on the generative law. On the discrete-informed axes the Laplace
posterior standard deviation runs slightly below the residual error (a light-tailed
approximation), so interval \emph{width} is marginally optimistic even though the
$95\%$-level coverage, fixed by the interval quantiles, is exact. The next subsection addresses
the conditioning limitation directly.

\subsection{Robustness of coverage under misspecification}
\label{sub:misspec}

The coverage study of \S\ref{sub:coverage} draws data from the fitted model, so it cannot
distinguish a well-calibrated estimator from one that is calibrated only because it scores data
it also generated. We therefore stress the projection with data drawn from processes it does
\emph{not} assume, and ask how far coverage drifts. Four one-parameter families of violations are
used, each targeting a specific assumption of \S\ref{sec:model}:
\begin{enumerate}\itemsep2pt
  \item \textbf{Correlated residuals} (violates local independence, \eqref{eq:localindep}): the
    residual of every item receives a shared nuisance component,
    $\varepsilon_j\leftarrow(\varepsilon_j+\rho\,g_\ast)/\sqrt{1+\rho^2}$ with common
    $g_\ast\sim\mathcal N(0,1)$ and $\rho\in\{0,0.15,0.30,0.45\}$.
  \item \textbf{Omitted factor} (violates the rank-$K$ assumption): a ninth latent
    $g_9\sim\mathcal N(0,1)$ loads a random $40\%$ of items with strength $s\in\{0,0.25,0.5,0.75\}$,
    absent from the scoring model.
  \item \textbf{Loading mismatch} (parameter error / transfer to a new sample): the scoring model
    uses $\lambda_{jk}(1+\varepsilon\xi_{jk})$, $\xi_{jk}\sim\mathcal N(0,1)$,
    $\varepsilon\in\{0,0.10,0.20,0.30\}$, while data use the true $\lambda_{jk}$.
  \item \textbf{Heavy tails} (violates Gaussian residuals): continuous residuals are Student-$t$
    with $\mathrm{df}\in\{\infty,8,5,3\}$, rescaled to unit variance.
\end{enumerate}
Each cell draws $N=4{,}000$ synthetic patients keeping $50\%$ of items (three seeds), scores them
with the \emph{unperturbed} Gaussian EAP, and records empirical coverage of the nominal $95\%$
per-axis interval. The results (\MFig{fig:blockers}B; \MSec{sec:res-utility-prognosis})
show graceful, monotone degradation and no collapse:
\begin{center}\small
\begin{tabular}{lcccc}
\toprule
Violation & none & low & med & high\\
\midrule
Correlated residuals ($\rho$) & $0.951$ & $0.943$ & $0.918$ & $0.889$\\
Omitted ninth factor ($s$) & $0.951$ & $0.948$ & $0.939$ & $0.939$\\
Loading mismatch ($\varepsilon$) & $0.949$ & $0.950$ & $0.927$ & $0.934$\\
Heavy-tailed residuals ($\mathrm{df}$) & $0.949$ & $0.950$ & $0.950$ & $0.955$\\
\bottomrule
\end{tabular}
\end{center}
Three features are worth noting. First, the unperturbed baseline recovers nominal ($0.949$--$0.951$)
in every family, confirming the estimator is unbiased when its assumptions hold. Second, heavy
tails are absorbed essentially without loss --- indeed slightly conservatively at the $50\%$
level --- because the EAP posterior variance depends on the loadings and the residual
\emph{variance}, not on higher residual moments. Third, the steepest degradation is under
correlated residuals, and this is expected: local dependence inflates the effective information a
naive independent-likelihood score believes it has, shrinking intervals below their honest width.
Even so, an extreme shared-nuisance loading ($\rho=0.45$, a residual correlation the
model-selection diagnostics of the main text would flag) costs only six points of coverage. The
projection's intervals are thus honest under the model and degrade predictably, not
catastrophically, away from it; local dependence is the assumption whose violation matters most,
and the one most worth checking on real data before trusting interval \emph{widths}.

%=============================================================================
\section*{Summary of cross-references to the main text}
\addcontentsline{toc}{section}{Cross-references to the main text}

Each numbered result here underlies a specific claim in the main manuscript:
\eqref{eq:obslik} and Proposition~\ref{prop:marg} justify the no-imputation invariant
(\MEq{eq:obslik}); Lemma~\ref{lem:woodbury} and Proposition~\ref{prop:eap} are the projection
(\MEq{eq:m-woodbury}, \MEq{eq:m-scores}); Propositions~\ref{prop:latslope}--\ref{prop:respslope}
are the loading-as-slope (\MEq{eq:m-slope}); Propositions~\ref{prop:additive}--\ref{prop:families}
and Corollary~\ref{cor:reliability} are the information accounting (\MEq{eq:m-precision},
\MEq{eq:m-fisher}); Propositions~\ref{prop:upperbound}--\ref{prop:voi} and
Corollary~\ref{cor:greedy} are adaptive sampling and value of information (\MEq{eq:m-voi});
Definition~\ref{def:r2}, Proposition~\ref{prop:kd} (the random-subspace $k/d$ bound that reframes
reconstruction $R^2$) and Proposition~\ref{prop:fog} are the archetype reconstruction; and
Propositions~\ref{prop:exact}--\ref{prop:priorout} with the misspecification stress test of
\S\ref{sub:misspec} are the calibration protocol and its robustness (\MSec{sec:m-calib},
\MSec{sec:res-utility-prognosis}).

\bibliographystyle{unsrtnat}
\bibliography{references}

\end{document}
